\Chapter{Wyniki i ograniczenia}
\label{ch:discussion}

Wynikiem bieżącego etapu jest działająca część infrastruktury danych
oraz sprawdzalna diagnoza ich ograniczeń. Eksport i niezależny od cache
weryfikator umożliwiają przekazanie zespołowi tej samej wersji wejścia.
Zachowanie surowych wierszy pozwala wrócić do niejednoznacznych wizyt
w alei, zamiast utrwalać niepewne odpowiedzi już podczas importu.
To warunek dalszego eksperymentu, ale nie dowód osiągnięcia celu
predykcyjnego.

Najważniejsze ograniczenie jest semantyczne. Poprawne typy i zgodne
skróty plików nie rozstrzygają, czy opisane zdarzenie było wymianą opon.
Audyt podważył wskaźnik oparty na \code{Stint}; wysoka zgodność z rejestrem
wizyt Jolpica nie usuwa problemu, ponieważ oba źródła mogą opisywać
przejazdy związane z karą. Wdrożenie etykiety przed rozstrzygnięciem D2
mogłoby doprowadzić do uczenia innego zadania niż zadeklarowane.

Drugie ograniczenie dotyczy dostępności czasowej. Zachowane zapowiedzi
liczby okrążeń umożliwią zbadanie, co było wiadomo w chwili $T$, ale
nie zastępują kontraktu ich użycia. Trzeba też obsłużyć początek historii
przed pierwszą znaną zapowiedzią. Podobnie poprawny kształt okna
sekwencyjnego nie dowodzi, że w jego obliczeniu nie wykorzystano
przyszłości innego kierowcy. Dopiero implementacja i testy przyczynowe
pozwolą taką własność potwierdzić.

Ograniczenie populacji do wyścigów suchych i bez czerwonej flagi zmniejsza
liczbę trudnych sytuacji, ale zawęża przyszłe zastosowanie modelu.
Nie będzie można bez osobnej oceny przenosić wniosków na deszcz,
przerwane wyścigi, sprinty ani następne sezony. Nawet w przyjętym
zakresie braki historii opon i szczególne regulacje wymagają jawnego
raportowania. Określenie kompletności w manifeście odnosi się do
żądanego eksportu i zapisanych kontroli, a nie do pełnej obserwowalności
procesu podejmowania decyzji przez zespoły.

Na tym etapie nie można odpowiedzieć, która metoda prognozuje najlepiej,
jaki jest poziom fałszywych alarmów ani czy sieci przewyższają LightGBM.
Nie ma podstaw do rysowania krzywych wyników ani wypełniania tabel
przykładowymi metrykami. Po ukończeniu pipeline'u użyteczność powinna
być oceniana na zamrożonym protokole, z podziałem na zieloną flagę,
SC/VSC i wyścigi ze szczególnymi regułami opon.

Realistyczny przyszły panel może wspierać analizę powtórek: pokazywać
przebieg prawdopodobieństwa, moment alarmu oraz zaobserwowaną wymianę
opon. Wynik pozostanie prognozą zachowania w warunkach danych
historycznych. Ocenę optymalnego momentu zjazdu należałoby oprzeć na
innym zadaniu, modelu skutków decyzji i dodatkowych dowodach.
